% 你或许需要的：

\BookChapter{ }{  }%%
\UnnumberedSection{章节备注}{Chapter Notes}
\begin{bilingualtext}
\bi{% 中

}{% 英

}
\end{bilingualtext}%%%%%%
\UnnumberedSection{正文}{Main Text}
\begin{bilingualtext}

\end{bilingualtext}%%%%%%



% \Bignote{正文}{Main Text}

\begin{center}
\raisebox{5mm}{\rule{0.6\linewidth}{0.4pt}}
\end{center}
\vspace{0.3pt}%%%%%%%%%%%

% 段落格式
\bi{% 中

}{% 英

}


\frontmatter

% 封面
\begin{titlepage}
    \thispagestyle{empty}
    \centering
    \vspace*{25mm}
    {\fontsize{30pt}{36pt}\selectfont 书名}
    \vspace{8mm}
    {\fontsize{13pt}{17pt}\selectfont Subtitle}
    \vfill
    {\fontsize{11pt}{15pt}\selectfont 作者姓名}
\end{titlepage}

\cleardoublepage

% 扉页
\thispagestyle{empty}
\begin{center}
    \vspace*{35mm}
    {\fontsize{25pt}{31pt}\selectfont 书名}
    \vspace{5mm}
    {\fontsize{12pt}{16pt}\selectfont Subtitle}
    \vfill
    作者姓名
\end{center}

\cleardoublepage

% 版权页
\thispagestyle{empty}
\vspace*{\fill}
{\small
\textbf{书名}
\vspace{4mm}

第一版
\vspace{4mm}

© 2026 作者姓名
\vspace{4mm}

All rights reserved.
\vspace{10mm}

本书内容仅供学习与研究使用。
}
\vspace*{25mm}

\cleardoublepage

% 目录
\tableofcontents

\mainmatter

% 序章
\BookChapter{序章}{Prologue}
这里是一段正文。
这是一条注释示例\footnote{这里就是注释。使用标准的 \texttt{\textbackslash footnote\{\}} 即可，LaTeX 会自动把注释放到当前页面底部。}。
这里继续正文。
\vspace{1em}
\bookdate{2026 年 9 月}

% 第一章
\BookChapter{第一章：开端}{Chapter One: The Beginning}
\ifbilingual
    \begin{bilingualtext}
        \bi{世界从一个极其微小的瞬间开始。}
           {The world begins with an extraordinarily small moment.}
        \bi{有些事情发生以前，谁也不知道它究竟意味着什么。}
           {Before something happens, no one knows what it will eventually mean.}
        \bi{我们只能在时间向前移动以后，重新回头理解它。}
           {We can only look back and understand it after time has moved forward.}
        \bi{这是一条带有注释的句子\footnote{双语模式下，脚注仍然按照正常的脚注机制处理，并排版在页面底部。}。}
           {This sentence contains a footnote\footnote{In bilingual mode, footnotes continue to use the standard bottom-of-page mechanism.}.}
    \end{bilingualtext}
\else
    世界从一个极其微小的瞬间开始。
    有些事情发生以前，谁也不知道它究竟意味着什么。
    我们只能在时间向前移动以后，重新回头理解它。
    这是一条带有注释的句子\footnote{这里是普通中文正文中的脚注示例。}。
\fi

% 第二章
\BookChapter{第二章：时间}{Chapter Two: Time}
\BookSection{过去}{The Past}
\ifbilingual
    \begin{bilingualtext}
        \bi{时间并不是一条笔直的线。}
           {Time is not a straight line.}
        \bi{当我们重新回忆过去时，过去也会被现在改变。}
           {When we remember the past, the past itself is changed by the present.}
    \end{bilingualtext}
\else
    时间并不是一条笔直的线。
    当我们重新回忆过去时，过去也会被现在改变。
\fi
\BookSection{现在}{The Present}
\ifbilingual
    \begin{bilingualtext}
        \bi{此刻是唯一真正发生着的事情。}
           {The present is the only thing that is truly happening.}
    \end{bilingualtext}
\else
    此刻是唯一真正发生着的事情。
\fi

% 后记
\backmatter
\chapter*{后记}
\addcontentsline{toc}{chapter}{后记}
这里是后记。